
\input header.tex
\footline={\hss\tenrm R-\folio\hss}\pageno=1

\def\3i{\indent\indent\indent}
\font\tencyr=wncyr10
\input cyracc.def
\def\cyr{\tencyr\cyracc}

\centerline{\bf REFERENCES$\qquad$}
\bigskip

\indent\indent\indent{\bf For general background, see for example:}\smallskip

\ref [A1] L. Ahlfors, Complex Analysis, McGraw-Hill 1966.

\ref [Mu] J. Munkres, Topology: A First Course, Prentice-Hall 1975.

\ref [Wi] T. Willmore, An Introduction to Differential Geometry, Clarendon, 1959.
\medskip

\3i{\bf Surveys of conformal dynamics.}
\smallskip
\ref [Bl1] P. Blanchard, Complex analytic dynamics on the Riemann sphere,
Bull. Amer. Math. Soc. {\bf 11} (1984), 85-141.

\ref [BlC] P. Blanchard and A. Chiu, Conformal dynamics: an informal
discussion, Lecture Notes, Boston University 1990.

\ref [Br] H. Brolin, Invariant sets under iteration of rational functions,
Arkiv f\"or Mat. {\bf 6} (1965) 103-144. % (See p. 118.)

\ref [C] L. Carleson, Complex Dynamics, Lecture Notes UCLA 1990.

\ref [Dv1] R. Devaney, An Introduction to Chaotic Dynamical Systems,
$2^{nd}$ ed., Addison-Wesley, 1989.

\ref [D1] A. Douady, Syst\`emes dynamiques holomorphes, S\'eminar Bourbaki, $35^e$
ann\'ee 1982\--83, ${\rm n}^o$ 599.

\ref [Fa] K. Falconer, Fractal Geometry: Mathematical Foundations and
Applications, Wiley 1990. (Ch. 14.)

\ref [K1] L. Keen, Julia sets, pp. 57-74 of ``Chaos and Fractals, the Mathematics
behind the Computer Graphics'', edit. Devaney and Keen, Proc. Symp. Appl.
Math. {\bf 39}, Amer. Math. Soc. 1989.

\ref [L1] M. Lyubich, The dynamics of rational transforms: the topological picture,
Russian Math. Surveys {\bf 41:4} (1986), 43-117.

\medskip
\3i{\bf \S1. Simply connected surfaces; uniformization.}\smallskip

\ref [A2] L. Ahlfors, Conformal Invariants, McGraw-Hill 1973.

\ref [AS] L. Ahlfors and L. Sario, Riemann Surfaces, Princeton U. Press 1960.

\ref [Be] A. Beardon, A  Primer on Riemann Surfaces, Cambridge U. Press 1984.

\ref [FK] H. Farkas and I. Kra, Riemann Surfaces, Springer 1980.

\ref [Sp] G. Springer, Introduction to Riemann Surfaces, Addison-Wesley 1957.

\medskip
\3i{\bf \S2. Montel's theorem.}\smallskip

\ref [Mo] P. Montel, Le\c cons sur les Familles Normales, Gauthier-Villars 1927.
\medskip

\3i{\bf \S3. Fatou and Julia.}\smallskip

\ref [Ca] A. Cayley, Application of the Newton-Fourier method to an
immaginary root of an equation, Quart. J. Pure Appl. Math. {\bf 16} (1879)
179-185.

\ref [F1] P. Fatou, Sur les solutions uniformes de certaines \'equations
fonctionnelle, C. R. Acad. Sci. Paris {\bf 143} (1906) 546-548.

\ref [F2] P. Fatou, Sur les \'equations fonctionnelles, Bull. Soc. math. France
{\bf 47} (1919) 161-271, and {\bf 48} (1920) 33-94, 208-314.

\ref [J] G. Julia, Memoire sur l'iteration des fonctions rationelles,
J. Math. Pure Appl. {\bf 8} (1918) 47-245.

\ref [Ri] J. F. Ritt, On the iteration of rational functions, Trans. Amer.
Math. Soc. {\bf 21} (1920) 348-356.
\medskip

\3i{\bf \S4. Iterated maps (Hyperbolic and Euclidean cases).}
\smallskip

\ref [Ba1]
I. N. Baker, Repulsive fixedpoints of entire functions, Math. Zeit. {\bf 104}
(1968) 252-256.

\ref [Ba2], ---------, An entire function which has wandering domains,
J. Austral. Math. Soc. {\bf 22} (1976) 173-176.

\ref [Dn1] A. Denjoy, Sur l'it\'eration des fonctions analytiques, C.R.Acad.Sci.
Paris {\bf 182} (1926) 255-257.

\ref [Dv2] R. Devaney, Exploding Julia sets, pp. 141-154 of ``Chaotic
Dynamics and Fractals'', edit. Barnsley and Demko, Academic Press 1986.

\ref [EL] A. Eremenko and M. Lyubich, Dynamical properties of some classes of
entire functions, Stony Brook Institute for Mathematical Sciences
Preprint 1990\#4. (See also Sov. Math. Dokl. {\bf 30} (1984) 592-594;
Func. Anal. Appl. {\bf 19} (1985) 323-324; and
J. Lond. Math. Soc. {\bf 36} (1987) 458-468.)

\ref [F3] P. Fatou, Sur l'it\'eration des fonctions transcendantes enti\`eres,
Acta Math. {\bf 47} (1926) 337-370.

\ref [GK] L. Goldberg and L. Keen, A finiteness theorem for a dynamical class
of entire functions, Erg. Th. \& Dy. Sy. {\bf 6} (1986) 183-192.

\ref [K2] L. Keen, The dynamics of holomorphic self-maps of \[\C^*\], in
``Holomorphic Functions and Moduli'', edit. Drasin et al., Springer 1988.

\ref [L2] M. Lyubich, The measurable dynamics of the exponential map,
Siber. J. Math. {\bf 28} (1987) 111-127. (See also Sov. Math. Dokl.
{\bf 35} (1987) 223-226.)

\ref [R1] M. Rees, The exponential map is not recurrent, Math. Zeit. {\bf
191} (1986) 593-598.
\medskip

\3i{\bf \S5. Smooth Julia sets.}\smallskip

\ref [La] S. Latt\`es, Sur l'iteration des substitutions rationelles et les
fonctions de \Poincare, C. R. Acad. Sci. Paris {\bf 16} (1918) 26-28.

\ref [He1] M. Herman, Exemples de fractions rationnelles ayant une orbite dense
sur la sph\`ere de Riemann, Bull. Soc. Math. France {\bf 112} (1984) 93-142.

\ref [R2] M. Rees, Ergodic rational maps with dense critical point forward
orbit, Erg. Th. \& Dy. Sy. {\bf 4} (1984) 311-322.

\ref [R3] M. Rees, Positive measure sets of ergodic rational maps, Ann. Sci.
\'Ecole Norm. Sup. (4) {\bf 19} (1986) 383-407.

\ref [UvN] S. Ulam and J. von Neumann, On combinations of stochastic and
deterministic processes, Bull. Amer. Math. Soc. {\bf 53} (1947) 1120.
\medskip
\eject
\3i{\bf \S6. Attracting and repelling fixed points.}\smallskip

\ref [Sch] E. Schr\"oder, Ueber iterirte Functionen, Math. Ann. {\bf 3} (1871);
see p. 303.

\ref [K\oe ] G. K\oe nigs, Recherches sur les integrals de certains equations
fonctionelles, Ann. Sci. \'Ec. Norm. Sup. ($3^e$ ser.)
{\bf 1} (1884) suppl\'em. 1-41.

\ref [B\"o] L. E. B\"ottcher (~{\cyr B\"etkher\cdprime}),
The principal laws of convergence of iterates and
their application to analysis (Russian), Izv. Kazan. Fiz.-Mat. Obshch.
{\bf 14} (1904) 155-234.

\medskip

\3i{\bf \S7. Parabolic fixed points.}\smallskip

\ref [Le] L. Leau, \'Etude sur les equations fonctionelles \`a une ou plusi\`ers
variables, Ann. Fac. Sci. Toulouse {\bf 11} (1897).

\ref [\'Ec] J. \'Ecalle, Th\'eorie it\'erative: introduction a la th\'eorie des
invariants holomorphes, J. Math. Pure Appl. {\bf 54} (1975) 183-258.

\ref [Ca] C. Camacho, On the local structure of conformal mappings and holomorphic
vector fields, Ast\'erisque {\bf 59-60} (1978) 83-94.

\ref [MT] J. Milnor and W. Thurston, Iterated maps of the interval,
pp. 465-563 of ``Dynamical Systems (Maryland 1986-87)'', edit. J.C.Alexander,
Lect. Notes Math. {\bf 1342}, Springer 1988.

\ref({\bit See also~} [B\"o, pp.201-234] {\bit and~} [F3].)
\medskip

\3i{\bf \S8. Cremer points and Siegel disks.}\smallskip

\ref [Pf] G. A. Pfeifer, On the conformal mapping of curvilinear angles; the
functional equation \[\phi[f(x)]=a_1\phi(x)\], Trans. A. M. S. {\bf 18} (1917)
185-198.

\ref [Cr1] H. Cremer, Zum Zentrumproblem, Math. Ann. {\bf 98} (1927) 151-163.

\ref [Cr2] ------------,  \"Uber die H\"aufigkeit der Nichtzentren, Math. Ann. {\bf
115} (1938) 573-580.

\ref [Si] C. L. Siegel, Iteration of analytic functions, Ann. of Math. {\bf 43}
(1942) 607-612.

\ref [HW] G. H. Hardy and E. M. Wright, An Introduction to the Theory of Numbers,
Clarendon Press 1938, 1945, 1954.

\ref [Br] A. D. Bryuno, Convergence of transformations of differential equations
to normal forms, Dokl. Akad. Nauk USSR {\bf 165} (1965) 987-989.

\ref [SiM] C. L. Siegel and J. Moser, Lectures on Celestial Mechanics, Springer
1971.

\ref [He2] M. Herman, Recent results and some open questions on Siegel's
linearization  theorem of germs of complex analytic diffeomorphisms of
\[C^n\] near a fixed point, pp. 138-198 of Proc \[8^{th}\] Int. Cong. Math.
Phys., World Sci. 1986.

\ref [D2] A. Douady, Disques de Siegel et anneaux de Herman, S\'em. Bourbaki
39$^e$ ann\'ee (1986-87) n$^o$ 677.

\ref [Y1] J.-C. Yoccoz, Lin\'earisation des germes de diff\'eomorphismes
holomorphes de \[(\C,0)\], C. R. Acad. Sci. Paris {\bf 306} (1988)
55-58.
\medskip\eject

\3i{\bf \S9. Holomorphic fixed point formula.}\smallskip

\ref [AB] M. Atiyah and R. Bott, A Lefschetz fixed point formula for elliptic
differential operators, Bull.A.M.S. {\bf 72} (1966) 245-250.
\medskip

\3i{\bf \S10. Most periodic orbits repel.}\smallskip

\ref [Shi] M. Shishikura, On the quasiconformal surgery of rational functions,
Ann. Sci. \'Ec. Norm. Sup. {\bf 20} (1987) 1-29.
\medskip

% {\bf \S11. Repelling cycles are dense in the Julia set.}\smallskip

\3i{\bf \S12. Herman rings}\smallskip

\ref [Ar] V. Arnold, Small denominators I, on the mappings of the circumference
into itself, Amer. Math. Soc. Transl. (2) {\bf 46} (1965) 213-284.

\ref [CL] E. Coddington and N. Levinson, Theory of Ordinary Differential
Equations, McGraw-Hill 1955.

\ref [Dn2] A. Denjoy, Sur les courbes d\'efinies par les \'equations
diff\'erentielles \`a la surface du tore, Journ. de Math. {\bf 11} (1932)
333-375.

\ref [He3] M. Herman, Sur la conjugation diff\'erentiables des diff\'eomorphismes
du cercle \`a les rotations, Pub. I.H.E.S. {\bf 49} (1979) 5-233.

\ref [dM] W. de Melo, Lectures on One-Dimensional Dynamics, 17$^o$ Col.
Brasil. Mat., IMPA 1990.

\ref [Y2] J.-C. Yoccoz, Conjugation diff\'erentiable des diff\'eomorphismes du
cercle dont le nombre de rotation v\'erifie une condition diophantienne,
Ann. Sci. E.N.S. Paris (4) {\bf 17} (1984) 333-359.
 
\ref ({\bit See also~} [He1], [D2], [Shi].)

\medskip

\3i{\bf \S13. Classification of Fatou components.}\smallskip

\ref [S1] D. Sullivan, Conformal dynamical systems, pp. 725-752 of ``Geometric
Dynamics'', edit. Palis, Lecture Notes Math. {\bf 1007} Springer 1983.

\ref [S2] D. Sullivan, Quasiconformal homeomorphisms and dynamics I, solution
of the Fatou-Julia problem on wandering domains, Ann. Math. {\bf 122}
(1985) 401-418.

\noindent{(\bit See also~} [DH1] ?? {\bit and~} [DH2] {\bit below.})
\medskip

\3i{\bf \S14. Sub-hyperbolic and hyperbolic maps.}\smallskip

\ref [Sm] S. Smale, Differentiable dynamical systems, Bull. Amer. Math. Soc. {\bf
73} (1967) 747-817.

\ref [MSS] R. Ma\~n\'e, P. Sad and D. Sullivan, On the dynamics of rational maps,
Ann. Sci. \'Ec. Norm. Sup. Paris (4) {\bf 16} (1983) 193-217.

\ref [L3] M. Lyubich, Some typical properties of the dynamics of rational
maps, Russian Math. Surveys {\bf 38} (1983) 154-155. (See also Sov. Math.
Dokl. {\bf 27} (1983) 22-25/

\ref [L4] M. Lyubich, An analysis of the stability of the dynamics of rational
functions, Selecta Math. Sovietica {\bf 9} (1990) 69-90. (Russian original
published in 1984.)

\ref [DH1] A. Douady and J. H. Hubbard, A proof of Thurston's topological
characterization of rational functions, preprint, Mittag-Leffler 1984.

\ref [Shu] M. Shub, Global Stability of Dynamical Systems, Springer 1987.
\smallskip
{\bf Basic reference for \S\S14 and 16-19:}

\ref [DH2] A. Douady and J. H. Hubbard, \'Etude dynamique des polyn\^omes
complexes I \& II, Publ. Math. Orsay (1984-85).

\medskip

\3i{\bf \S15. Prime ends.}\smallskip

\ref [Ca] C. Carath\'eodory, \"Uber die Begrenzung einfach zusammenh\"angender
Gebiete, Math. Ann. {\bf 73} (1913) 323-370. (Gesam. Math. Schr., v. 4.)

\ref [Ep] D.B.A.Epstein, Prime Ends, preprint, Univ. Warwick 1978.

\ref [Ho] K. Hoffman, Banach Spaces of Analytic Functions, Prentice-Hall 1962.

\ref [Ma] J. Mather, Topological proofs of some purely topological consequences
of Carath\'eo\-dory's theory of prime ends, pp. 225-255 of ``Selected Studies'',
edit. T and G. Rassias, North-Holland 1982.

\ref [Oh] M. Ohtsuka, Dirichlet Problem, Extremal Length and Prime Ends, van
Nostrand 1970.

\noindent({\bit See also~} [A2].)
\medskip

\3i{\bf \S16. Local connectivity.}\smallskip

% \ref (The next few sections are based on the Douady-Hubbard Orsay Notes.
% See \S13 above.)\smallskip

\ref [Ku] K. Kuratowski, Topologie, Warsaw 1958-61.

\ref [HY] Hocking and Young, Topology, Addison-Wesley 1961.

\noindent({\bit See also~} [Mu].)
\medskip

\3i{\bf \S17. The filled Julia set.}\smallskip

\ref [Bro] M. Brown, A proof of the generalized Schoenflies theorem, Bulletin
A.M.S. {\bf 66} (1960) 74-76. (See also ``The monotone union of open
n-cells is an open n-cell'', Proc.A.M.S. {\bf 12} (1961) 812-814.)

\ref [Bl2] P. Blanchard, Disconnected Julia sets, pp. 181-201 of ``Chaotic
Dynamics and\break Fractals'', edit. Barnsley and Demko, Academic Press 1986.

\ref [BH1] B. Branner and J. H. Hubbard, The iteration of cubic polynomials,
Part I: the global topology of parameter space, Acta Math. {\bf 160} (1988)
143-206.

\ref [BH2] B. Branner and J. H. Hubbard, The iteration of cubic polynomials,
Part II: patterns and parapatterns, Acta Math., to appear.\medskip


\3i{\bf \S18. External rays and periodic points.}\smallskip

\ref [Pe] C. Petersen, On the Pommerenke-Levin-Yoccoz inequality,
preprint, IHES 1991.

\ref [Po] C. Pommerenke, On conformal mapping and iteration of rational
functions, Complex Var. Th. \& Appl. {\bf 5}, 1986.

\ref [G] L. Goldberg, Rotation cycles on the unit circle, Stony Brook
IMS preprint 1990/14.

\ref [GM] L. Goldberg and J. Milnor, Fixed point portraits of polynomial maps,
Stony Brook IMS preprint 1990/14.\medskip
% \eject

\3i{\bf Appendix A. Theorems from classical analysis.}\smallskip

\ref [Je] J. L. W. V. Jensen, Sur un nouvel et important th\'eoreme de la th\'eorie
des fonctions, Acta Math. {\bf 22} (1899) 219-251.

\ref [RR] F. and M. Riesz, \"Uber Randwerte einer analytischen Funktionen,
Quatr. Congr. Math. Scand. Stockholm 1916, 27-44.

\ref [Ko] P. Koebe, \"Uber die Uniformizierung beliebiger analytischer Kurven,
Nachr. Akad. Wiss. G\"ottingen, Math.-Phys. Kl. (1907) 191-210.

\ref [Gr] T. H. Gronwall, Some remarks on conformal representation, Ann. of
Math. {\bf 16} (1914-15) 72-76.

\ref [Bi] L. Bieberbach, \"Uber die Koeffizienten derjenigen Potenzreihen
welche eine schlichte Abbildung des Einheitskreises vermitteln, S.-B.
Preuss. Akad. Wiss. (1916) 940-955.

\ref [dB] L. de Brange, A proof of the Bieberbach conjecture, Acta Math. {\bf
154} (1985) 137-152.
\medskip

\3i{\bf Appendix B. Length-area-modulus inequalities.}\smallskip

\ref [CR] R. Courant and H. Robbins, What is Mathematics?, Oxford U. Press
 1941.\smallskip

\ref ({\bit See also~} [A2], [BH2].)\medskip

\3i{\bf Appendix C. Continued fractions.}\smallskip

\ref ({\bit See~} Hardy and Wright [HW].)\medskip

\3i{\bf Appendix D. Two complex variables.}\smallskip

\ref [Ze] E. Zehnder, A simple proof of a generalization of a theorem by C. L.
Siegel, in ``Geometry and Topology III'', edit. do Carmo and Palis, Lecture
Notes Math. {\bf 597}, Springer 1977.

\ref [Ne] S. Newhouse, Diffeomorphisms with infinitely many sinks, Topology
{\bf 16} (1974) 9-18.

\ref [H] J. Hubbard, The H\'enon mapping in the complex domain, pp. 101-111 of
Chaotic Dynamics and Fractals, M. Barnsley and S. Demko, (ed.),
Academic Press 1986.

\ref [FM] S. Friedland and J. Milnor, Dynamical properties of plane polynomial
automorphisms, Erg. Th. \& Dy. Sy. {\bf 9} (1989) 67-99.

\ref [Be] E. Bedford, Iteration of polynomial automorphisms of \[\C^2\],
Preprint, Purdue 1990.

\medskip

\3i{\bf Appendix E. Branched coverings and orbifolds}\smallskip

\ref [Mas] B. Maskit, Kleinian Groups, Grundl. math. Wiss.  {\bf 287} Springer
1987.

\ref [M1] J. Milnor, On the 3-dimensional Brieskorn manifolds \[M(p,q,r)\],
pp. 175-225 of ``Knots, Groups, and 3-Manifolds'', edit. Neuwirth, Ann.
Math. Studies {\bf 84}, Princeton U. Press 1975. (See \S2.)

\ref [T] W. Thurston, Three-dimensional Geometry
and Topology, to appear.
\medskip\eject

\medskip\3i{\bf Appendix F. Parameter space.}\smallskip

\ref [BM] R. Brooks and P. Matelski, The dynamics of 2-generator subgroups
of \[PSL(2,\C)\], pp. 65-71 of ``Riemann Surfaces and Related Topics'',
Proceedings 1978 Stony Brook Conference, edit. Kra and Maskit,
Ann. Math. Stud. {\bf 97} Princeton U. Press 1981.

\ref [Man] B. Mandelbrot, Fractal aspects of the iteration of \[z\mapsto\lambda
z(1-z)\] for complex \[\lambda,\,z\], Annals NY Acad. Sci. {\bf 357} (1980)
249-259.

\ref [DH3] A. Douady and J. H. Hubbard, It\'eration des polyn\^omes quadratiques
complexes, CRAS Paris {\bf 294} (1982) 123-126.

\ref [DH4] A. Douady and J. H. Hubbard, On the dynamics of polynomial-like
mappings, Ann. Sci. Ec. Norm. Sup. (Paris) {\bf 18} (1985), 287-343.

\ref [D3] A. Douady, Julia sets and the Mandelbrot  set, pp. 161-173
of ``The Beauty of Fractals'', edit. Peitgen and Richter, Springer 1986.

\ref [B1] B. Branner, The Mandelbrot set, pp. 75-105
of ``Chaos and Fractals'', edit.
Devaney and Keen, Proc. Symp. Applied Math. {\bf 39}, Amer. Math. Soc. 1989.

\ref [B2] B. Branner, The parameter space for complex cubic polynomials,
pp. 169-179 of ``Chaotic Dynamics and fractals'', edit. Barnsley and Demko,
Academic Press 1986.

%\ref [BH2] B. Branner and J. H. Hubbard, The iteration of cubic polynomials,
%Part I: the global topology of parameter space, Acta Math. {\bf 160} (1988)
%143-206.

\ref [M2] J. Milnor, { Remarks on iterated cubic maps}, Stony Brook IMS
preprint 1990/6.

\ref [M3] J. Milnor, { Hyperbolic components in spaces of polynomial maps},
in preparation.

\ref [M4] J. Milnor, { On cubic polynomials with periodic critical point},
in preparation.

\ref [R4] M. Rees, Components of degree two hyperbolic rational maps,
Invent. Math. {\bf 100} 1990, 357-382.

\ref [R5] M. Rees, { A partial description of parameter space of rational
maps of degree two}, Part I: preprint, Univ. Liverpool 1990;
and Part II: Stony Brook IMS preprints 1991/4.

\ref({\bit See also~} [BH1].)\medskip

\3i{\bf Appendix G. Computer graphics.}\smallskip

\ref [M5] J. Milnor, Self-similarity and hairiness in the Mandelbrot set,
pp. 211-257 of ``Computers in Geometry and Topology'', edit. Tangora,
Lect. Notes Pure Appl. Math. {\bf 114}, Dekker 1989 (cf. p. 218 and \S5).

\ref [Fi] Y. Fisher, Exploring the Mandelbrot set, pp. 287-296 of
The Science of Fractal Images,  edit. Peitgen and Saupe,
Springer 1989.

\ref [Pei] H.-O. Peitgen, Fantastic deterministic fractals, ibid. pp. 169-218.
\medskip


\vfill\bigskip

%\rightline{jwm, version of 8-30-91}
\eject

\footline={\hss\tenrm I-\folio\hss}\pageno=1

\centerline{\bf INDEX}\smallskip
$$\vbox{\settabs 2 \columns
\+Abel & 7-4\cr
\+annulus & 2-2, 2-5, 4-2, 12-1\cr
\+Arnold & 12-1, 12-2 \cr
\+attracting & 4-1, 3-2, 6-1, 7-1\cr
\+Baire's Theorem & 3-7, 8-1 \cr
\+ Baker & 4-1, 11-5\cr
\+basin & 4-2, 3-2, 7-1, 10-1\cr
\+Bieberbach & A-3ff, D-2\cr
\+Blaschke product & 5-3, 12-2, 12-4\cr
\+B\"otk$\!$her (= B\"ottcher) & 6-6, 17-1\cr
\+Branner & 17-3, F-2\cr
\+Brooks \& Matelski & F-1-F-2\cr
\+Brown & 17-2\cr
\+Bryuno & 8-5ff, 12-4\cr
\+Carath\'eodory & 1-1, \S15, \S16, 17-3, 18-1\cr
\+cellular set & 17-2, F-2\cr
\+complete surface & 1-4\cr
\+conformal automorphism, metric & 1-1,1-3\cr
\+continued fraction & 8-5, 12-4, \S C\cr
\+covering surface & 2-1, 2-5, E-1\cr
\+Cremer & \S8, 11-4, 18-3\cr
\+curvature & 1-3, 1-5, 1-9, 2-5\cr
\+cycle & 3-2, 10-1\cr
\+deck transformation & 2-1, E-1\cr
\+Denjoy & 4-1, 12-1\cr
\+Diophantine number & 8-3, 12-2 \cr
\+Douady & 14-1, 17-3, \S18, F-2, G-1\cr
\+\'Ecalle & 7-4\cr
\+Euclidean surface & 2-1\cr
\+Euler characteristic & 5-2, 14-4, E-2ff\cr
\+expanding map & 14-1\cr
\+external ray & 17-2, \S18\cr
\+Fatou & \S3, 4-1, 6-4, \S7, \S10, 11-2, 13-1, 15-3, D-2\cr
\+Flower Theorem & 7-1\cr
\+fractal & 5-1, 6-2, F-2\cr
\+fractional linear transformation & 1-2, 1-5, 1-9\cr
\+fully invariant & 3-2\cr
\+generic & 3-7, 8-1\cr
\+grand orbit & 3-2, 3-6\cr
\+Green's function & 17-1, G-2\cr
}$$ 
\vfil\eject

$$\vbox{\settabs 2 \columns
\+half-plane & 1-2\cr
\+H\'enon map & D-1\cr
\+Herman & 5-3, 8-8, \S12, 13-1\cr
\+homoclinic point & 11-1  \cr
\+Hubbard & 14-1, 17-3, \S18, \S F-1, G-1\cr
\+hyperbolic, sub-hyperbolic map & 14-1, 14-3, 17-3, F-1\cr
\+Hyperbolic surface & 2-2\cr
\+Jensen & A-1\cr
\+Jordan curve & 3-3, 6-4, 7-7, 15-4, 16-2\cr
\+Julia & \S3, 6-4, 11-1, 13-1\cr
\+Koebe & 1-1, A-4\cr
\+K\oe nigs & 6-1, D-1\cr
\+Latt\`es & 5-1ff\cr
\+Leau & \S7\cr
\+Lebesgue measure & 8-2, 8-8, 15-2, 15-3, A-2\cr
\+Liouville & 1-2, 8-3, 8-10\cr
\+local connectivity & \S16, 17-3, 18-1ff, F-2\cr
\+locally uniform convergence & 2-3\cr
\+Lyubich & 4-1, 8-2, F-2\cr
\+Mandelbrot & \S F, \S G \cr
\+maximum modulus principle & 1-1, 17-1\cr
\+Montel's Theorem& 2-4, 3-6, 11-2\cr
\+multiplier & 3-2\cr
\+neutral & 3-2\cr
\+Newton's method & 5-3, F-1\cr
\+normal family & 2-4\cr
\+$o(),\;{\cal O}()$ & 6-6, 7-3\cr
\+orbifold, orbifold metric & 14-3ff, \S E\cr
\+parabolic & 1-7, 1-9, 2-2, \S7\cr
\+periodic orbit & 3-2, \S10, \S11\cr
\+petal & 7-1, 7-8, 18-6\cr
\+Pfeiffer & 8-1\cr
\+Picard's Theorem & 2-3\cr
\+\Poincare  & 1-3ff, 2-2, 12-1\cr
\+prime end & 15-5\cr
\+proper map & 5-2, 12-4, 13-1, E-1\cr
\+PSL & 1-2, 1-6, 1-9\cr
\+punctured plane, disk & 2-1, 2-2, 4-2\cr
\+quasi-conformal deformation, surgery & 10-1, 12-1, 12-4, 13-4\cr
}$$
\vfil\eject

$$\vbox{\settabs 2 \columns
\+ramification & 14-4, E-1\cr
\+rational function & 3-1, 5-3\cr
\+Rees & 4-1, 5-3, 8-2, F-2\cr
\+repelling & 3-2, 6-1, 7-1\cr
\+residue & 9-1 \cr
\+Riemann surface, metric, sphere, map& 1-1, 1-3, 1-5, 15-1\cr
\+Riemann-Hurwitz &5-2, 17-2, E-2 \cr
\+Riesz & 8-9, 15-3, 16-4, A-2\cr
\+root of unity & 7-1, 7-4\cr
\+rotation, rotation number & 1-2, 4-1ff, 12-1ff, 13-1\cr
\+Schr\"oder & 6-1\cr
\+Schwarz Lemma, & 1-1ff, A-4\cr
\+Schwarz inequality&  15-2, B-1\cr
\+self-similarity & 3-2, 3-9\cr
\+Shishikura & 10-1, 12-1, 12-4\cr
\+Siegel & 8-2ff, 11-3, 12-1\cr
\+size (of disk) & 8-8\cr
\+Smale & 14-3\cr
\+Snail Lemma & 13-2, 18-5 \cr
\+spherical metric, surface & 1-5, 2-1\cr
\+spiral & 6-2\cr
\+Sullivan & \S13-1, 14-3, 18-2, F-2\cr
\+superattracting & 3-2, 6-6\cr
\+Tchebycheff polynomial & 5-3\cr
\+thrice punctured sphere & 2-3\cr
\+Thurston & 11-4, 13-4, 14-3, \S E \cr
\+torus & 2-1, 4-1, 4-6, 5-1ff\cr
\+transcendental function & 4-1\cr
\+transitive action & 1-2, 1-8\cr
\+Ulam \& Von Neumann & 5-1\cr
\+uniformization & 1-1, 2-1\cr
\+univalent & 7-4, 15-3, A-1ff\cr
\+wandering component, point & 13-4, 14-3\cr
\+Weierstrass Theorem, \[\wp\] & 1-2, 5-2\cr
\+Yoccoz & 8-6ff, 12-2, 18-2, 18-6, F-2\cr
}$$

\end



























